\documentclass[12pt, a4paper]{article}

\usepackage[margin=2.5cm]{geometry}
\usepackage{hyperref}
\usepackage{booktabs}
\usepackage{array}
\usepackage{parskip}
\usepackage{titlesec}
\usepackage{xcolor}
\usepackage[T1]{fontenc}
\usepackage[utf8]{inputenc}

\hypersetup{
    colorlinks=true,
    linkcolor=blue!60!black,
    urlcolor=blue!60!black,
    citecolor=blue!60!black
}

\titleformat{\section}{\large\bfseries}{}{0em}{}[\titlerule]
\titleformat{\subsection}{\normalsize\bfseries}{}{0em}{}

\title{\textbf{Seminar Thesis on \textit{Generating Computational Cognitive Models
using Large Language Models}}\\
\large Submitted by Christian Block and YingZhu Chen}

\author{}
\date{}

\begin{document}

\maketitle
\tableofcontents
% ─────────────────────────────────────────────────────────────
\section{Executive Summary}

The pipeline proposed in  \textit{Generating Computational Cognitive Models
using Large Language Models} (GeCCo) uses Large Language Models (LLMs) to automatically generate, fit, and iteratively refine computational cognitive models as Python functions. Given a task description, participant behavioural data, and a template function, the pipeline in GeCCo prompts an LLM to propose candidate models, fits those proposals to held-out data, and refines them based on feedback derived from predictive performance.

The approach was benchmarked across four cognitive domains---decision making, learning, planning, and memory---using three open-source LLMs of varying parameter size, capacity, and training procedure. On all four human behavioural datasets, LLM-generated models consistently matched or outperformed the best domain-specific models from the cognitive-science literature, while remaining interpretable.

% ─────────────────────────────────────────────────────────────
\section{Motivation}

Computational cognitive models provide a formal framework for understanding human behaviour. In practice, researchers hand-craft these models, fit them to behavioural data, and iteratively refine them. This process is slow, skill-intensive, and systematically biased due to person's past exposure to the cognitive literature. Translating a psychological theory into a working algorithm requires lengthy literature reviews and repeated coding cycles. Researchers must simultaneously act as domain experts in cognitive science, theoreticians for data modeling, and software engineers for implementation. Perhaps most importantly, researchers tend to explore only the model classes they already know and favour, potentially missing alternative explanations for the data. The  aforementioned limitations limit the space of models considered, motivating an automated approach that is not biased towards any particular school of thought in theoretical cognitive scienece.\\

Three specific LLM capabilities make them well-suited for propsoing cognitive models for experimental data.\footnote{ The models were used: LLMs: Llama-3.1-Instruct-70B (Llama; Meta Plat-
forms 2024, DeepSeek-R1-Distill-Llama-3.1-70B (R1; Guo et al. 2025), and Qwen2.5-72B-Instruct
(Qwen; Yang et al. 2024). } First, LLMs can process task descriptions and behavioural data written in natural language, which make it highly adaptable across different experimental paradigms without the need to provide custom interface. Second, drawing on broad pre-training knowledge, they can identify structure in complex behavioural data and generate plausible psychological hypotheses derived from the training data. Third, LLMs can translate those hypotheses directly into executable Python functions. These three abilities have already been utilitzed in automated statistical modelling through probabilistic program generation, and GeCCo applies the same logic to cognitive science.

% ─────────────────────────────────────────────────────────────
\section{ GeCCo's Proposed Pipeline}

Each pipeline run consists of 10 sampling iterations, repeated across 5 independent runs per beforementioned cognitive domain. The process begins by providing the base LLM with a fixed prompt containing a natural language task description, a subset of behavioral data, strict instructions, a domain-specific code template, and starting from the second iteration onwards, iterative feedback. Using this information, the LLM generates three unique cognitive models, ensuring that each features a non-overlapping set of parameters. \\

Once generated, these models are parsed from the response, converted into executable Python functions, and have their specific parameter names extracted. They are then fitted to a held-out validation dataset that was deliberately excluded from the initial prompt. To prevent the models from getting stuck in local minima, this fitting process is initialized from 20 random starting points.\\

After the fitting stage, the models are evaluated using the Bayesian Information Criterion (BIC) to identify the single best-performing candidate among the three. Finally, the code for this winning model, along with a list of previously used parameter names, is gathered to act as feedback. This feedback is then injected into the prompt for the next iteration, effectively guiding the LLM's future proposals and preventing it from duplicating its past ideas

\section{Experiments}

Two complementary metrics are used throughout the evaluation. BIC balances goodness of fit against model complexity by penalising the number of free parameters; when two models fit equally well. The Exceedance Probability (EXP) is used in the final evaluation stage and quantifies the posterior probability that a given model provides the most prevalent explanation of observed behaviour across the population, translating raw BIC scores into an interpretable confidence measure.

\subsection{Experiment 1: Decision Making}

People were asked to choose between two fake products (Option A and Option B). They were given ratings from four different "experts," but they were also told that some experts were more trustworthy (had higher "validity") than others. The baseline was the probabilistic Weighted Additive Decision  model, which uses four separate parameters to weight the four features. The Llama-generated model achieved a substantially lower BIC (39.40 vs.\ 51.97; $p < .001$). Rather than requiring four parameters, the LLM reduced the problem to a single \emph{discount factor} that smoothly interpolates between three classic decision strategies---\emph{Take the Best}, \emph{Equal Weighting}, and \emph{Weighted Additive}---mirroring complex Bayesian inference but arriving at the solution naturally.

\subsection{Experiment 2: Learning}

Participants played a "slot machine" style game where they repeatedly chose between two options to win money. In a special version of this game, they didn't just see what they won, but they also saw what the other option would have paid out if they had chosen it. The baseline was a four-learning-rate variant of the Rescorla--Wagner model augmented by a perseveration parameter capturing choice stickiness. The LLM-generated model (BIC 77.97 vs.\ 85.24; EXP = 0.97) replaced the asymmetric learning rates with a clean separation between actual and counterfactual value traces, and replaced the perseveration parameter with a psychologically motivated \emph{forgetting parameter} that decays unchosen-option values over time, mirroring the limits of working memory.

\subsection{Experiment 3: Planning}

This is a two-step puzzle game. Participants pick a "magic carpet," which usually takes them to a specific mountain. Once at the mountain, they ask a genie for gold. However, sometimes a "strong wind" blows the carpet to the wrong mountain. This game tests whether people just blindly repeat actions that win them gold, or if they actually "plan ahead" by understanding the rules of the wind and the carpets. The baseline was a Hybrid model combining model-free and model-based value iteration, mixed by a free weighting parameter. The R1-generated model achieved a lower mean BIC (432.47 vs.\ 437.38; EXP = 0.85), although the BIC difference did not reach conventional significance ($p = .27$). By unifying the model-free/model-based dichotomy into a single system using two transition-dependent learning rates, and by omitting the temporal discounting parameter, the model aligned with recent theoretical calls to move beyond strict dual-system accounts of planning.

\subsection{Experiment 4: Working Memory}

People had to learn through trial-and-error which button to press when shown different shapes on a screen. Sometimes they only had to juggle 3 shapes in their head (low cognitive load), and sometimes they had to juggle 6 shapes (high cognitive load). The baseline was the RL-WM model (Collins \& Frank, 2012), in which a slow RL module and a fast but capacity-limited WM module operate independently, with a load-dependent mixing weight. The Llama-generated model outperformed the baseline substantially (BIC 267.25 vs.\ 275.20; EXP = 0.99). Contrary to the standard assumption that reinforcement learning is immune to cognitive load, the LLM's model assigned separate RL learning rates for low- and high-load conditions (\texttt{lr\_low} / \texttt{lr\_high}), discovering that load modulates the core learning signal---a conclusion consistent with neuroimaging evidence on reward prediction error signals.

\section{Control Experiments}


A mixed-effects regression predicting BIC from \emph{reasoning ability}, \emph{model family}, and \emph{model size} found that Llama outperformed Qwen ($\beta = 5.624$, $p = .03$) despite Qwen having more parameters, suggesting that training data and protocol matter more than raw scale. Reasoning ability showed only a marginal advantage ($\beta = -0.266$, $p = .07$), falling short of conventional significance.


Systematically removing each prompt component and refitting a mixed-effects model on post-ablation BIC scores revealed that all four components contribute meaningfully, though with different magnitudes. Iterative feedback was the strongest single driver of performance ($\beta = -17.040$, $p < .05$), followed by participant data ($\beta = -12.836$, $p < .05$), task description ($\beta = -9.834$, $p < .05$), and the code template ($\beta = -9.359$, $p < .05$). The dominance of the feedback component confirms that the iterative refinement loop---not merely the quality of the initial prompt---is what separates GeCCo from a single-shot query.


The LogProber method was applied to Llama prompts across all four domains. Acceleration scores were well below the contamination threshold of 1.0 (range: 0.017--0.648), confirming that the LLM is reasoning about novel problems rather than recalling memorised solutions from pretraining.


On synthetic datasets with known generating processes, the LLM successfully reverse-engineered the true model in both learning (Rescorla--Wagner) and decision-making (Take the Best; Tallying) domains, validating the mathematical soundness of the approach.

CENTAUR is a large foundation model trained to predict human behaviour across cognitive tasks, serving as a proxy for the ceiling of explainable variance. GeCCo matched or exceeded CENTAUR's negative log-likelihood across all four domains, while producing interpretable Python code rather than opaque model weights.

% ─────────────────────────────────────────────────────────────
\section{Discussion}
The GeCCo pipeline successfully demonstrates that open-source LLMs can act as automated research assistants, generating and refining psychological models using standard Python code. Across four distinct areas of human cognition, these AI-generated models performed just as well as, if not better than, the best hand-crafted models built by human experts. Rigorous control experiments proved this success was driven by the AI's ability to learn from its mistakes through iterative feedback, as well as its foundational reasoning architecture, rather than from data contamination or sheer luck. Furthermore, the AI was able to uncover the true underlying rules in simulated tests, and its resulting code remained highly interpretable to humans while capturing just as much predictable human behaviour as massive, opaque foundation models like CENTAUR.\\

A major benefit of GeCCo is its high accessibility; the system works out of the box using standard, open-source AI models without the need for expensive or time-consuming retraining. It utilizes a hybrid approach where the AI acts as an intelligent proposal engine to brainstorm the mathematical structures, but the actual testing of how well those models fit the data is done offline using standard, rigorous statistical tools. This ensures that the final model selection is anchored strictly in hard data, rather than relying on unverified AI outputs. Because it can seamlessly adapt to different types of psychological tasks, it provides scientists with a highly flexible tool to rapidly explore a much broader landscape of scientific theories.\\

Despite these breakthroughs, the authors outline several limitations that chart the course for future research. One limitation is template bias; the AI is currently provided with a basic, human-made code template to guide its syntax, which might unintentionally restrict its creative exploration to specific model classes. However, ablation tests showed this template had a minimal impact on performance, and allowing the AI to generate its own starting template proved to be a successful workaround. Another limitation involves the system's simplistic feedback mechanism, which currently only feeds the AI numerical scores, such as the Bayesian Information Criterion, to communicate how well its previous guess performed. Future iterations could incorporate a secondary AI to provide qualitative feedback---critiquing the theoretical logic, simplicity, or code quality---which could further minimize the need for human oversight. 

Additionally, the study is limited in scope because it evaluated four specific task domains, which do not represent the full spectrum of cognitive science. Future research must extend this pipeline to more complex areas, such as visual perception and language processing, which involve much larger datasets and require models with significantly more parameters. Finally, to truly prove its broader ecological validity beyond controlled laboratory settings, the pipeline needs to be evaluated using naturalistic datasets that capture the messy, unpredictable nature of human behaviour in the real world. Ultimately, these findings suggest that LLMs have the potential to democratize the highly technical field of computational modeling and significantly accelerate the pace of scientific discovery.
\subsection*{Context Within Related Work}
This iterative optimization paradigm closely mirrors the broader industry shift toward recursive self-improvement models and automated research engines, such as Sakana AI's ``The AI Scientist'', Google DeepMind's autonomous research loops for scientific discovery, and Anthropic's research into iterative code-refinement. These modern frameworks employ an AI system to continuously write, test, critique, and upgrade its own code to automate the scientific pipeline. GeCCo successfully adapts this concept to the behavioral sciences by turning the LLM into a rapid hypothesis proposal engine[cite: 410]. Crucially, whereas generalized self-improvement models risk hallucination or runaway optimization errors when critiquing themselves, GeCCo introduces a vital safeguard: the self-improvement cycle is anchored in a hybrid loop, where standard, rigid statistical frameworks provide the objective ground-truth feedback that drives the AI's refinement process.

\subsection{Connection to other Papers in the Seminar (TBC)}

% ─────────────────────────────────────────────────────────────
\section*{Reference}

Rmus et al.\ (2024). \textit{GeCCo: Guided Generation of Computational Cognitive Models}. Available at: \url{https://github.com/MilenaCCNlab/gecco.git}

\end{document}